% Copyright 2004 by Till Tantau <tantau@users.sourceforge.net>.
%
% Modified (c) 2022 by Paul Nong-Laolam
% for ESPEC North America, INC. Web Controller V3 training distribution 
%
% In principle, this file can be redistributed and/or modified under
% the terms of the GNU Public License, version 2.
%
% However, this file is supposed to be a template to be modified
% for your own needs. For this reason, if you use this file as a
% template and not specifically distribute it as part of a another
% package/program, I grant the extra permission to freely copy and
% modify this file as you see fit and even to delete this copyright
% notice. 

\documentclass{beamer}
% Replace the \documentclass declaration above
% with the following two lines to typeset your 
% lecture notes as a handout:
%\documentclass{article}
%\usepackage{beamerarticle}


% There are many different themes available for Beamer. A comprehensive
% list with examples is given here:
% http://deic.uab.es/~iblanes/beamer_gallery/index_by_theme.html
% You can uncomment the themes below if you would like to use a different
% one:
%\usetheme{AnnArbor}
%\usetheme{Antibes}
%\usetheme{Bergen}
%cis275-ch6-w22.pdf\usetheme{Berkeley}
%\usetheme{Berlin}
%\usetheme{Boadilla}
%\usetheme{boxes}
%\usetheme{CambridgeUS}
%\usetheme{Copenhagen}
%\usetheme{Darmstadt}
%\usetheme{default}
%\usetheme{Frankfurt}
\usetheme{Goettingen}  % selected for usage
%\usetheme{Hannover}
%\usetheme{Ilmenau}
%\usetheme{JuanLesPins}
%\usetheme{Luebeck}
%\usetheme{Madrid}
%\usetheme{Malmoe}
%\usetheme{Marburg}
%\usetheme{Montpellier}
%\usetheme{PaloAlto}
%\usetheme{Pittsburgh}
%\usetheme{Rochester}
%\usetheme{Singapore}
%\usetheme{Szeged}
%\usetheme{Warsaw}

% Figures
\usepackage{graphicx}
\usepackage{colortbl}
\usepackage{verbatim} 

% explore
\usepackage{wrapfig}    % wrap text aroudn figures 
\usepackage{caption}    % set control caption for minipage env
% end explore

% menu keys
\usepackage[os=win]{menukeys}        % use key obxes
\renewmenumacro{\keys}[+]{shadowedroundedkeys}

\title{ESPEC Web Controller V3}

% A subtitle is optional and this may be deleted
\subtitle{HALT/HASS T-Series Chambers}

\author{Paul Nong-Laolam}%\inst{1} } %\and S.~Another\inst{2}}
% - Give the names in the same order as the appear in the paper.
% - Use the \inst{?} command only if the authors have different
%   affiliation.

\institute[Universities of Somewhere and Elsewhere] % (optional, but mostly needed)
{
  %\inst{1}%
  ESPEC North America, Inc. \\ 
  4141 Central Parkway \\
  Hudsonville, Michigan 49426 (USA) 
  %\and
  %\inst{2}%
  %Department of Theoretical Philosophy\\
  %University of Elsewhere
  }
% - Use the \inst command only if there are several affiliations.
% - Keep it simple, no one is interested in your street address.
\date{1/25/2023}
% - Either use conference name or its abbreviation.
% - Not really informative to the audience, more for people (including
%   yourself) who are reading the slides online

\subject{Software} % \subjectComputer Science}
% This is only inserted into the PDF information catalog. Can be left
% out. 

% If you have a file called "university-logo-filename.xxx", where xxx
% is a graphic format that can be processed by latex or pdflatex,
% resp., then you can add a logo as follows:

% \pgfdeclareimage[height=0.5cm]{university-logo}{university-logo-filename}
% \logo{\pgfuseimage{university-logo}}

% Delete this, if you do not want the table of contents to pop up at
% the beginning of each subsection:
\AtBeginSubsection[]
{
  \begin{frame}<beamer>{Outline}
    \tableofcontents[currentsection,currentsubsection]
  \end{frame}
}

% PLACE PAGE NUMBERING
\addtobeamertemplate{navigation symbols}{}{%
    \usebeamerfont{footline}%
    \usebeamercolor[fg]{footline}%
    \hspace{1em}%
    \insertframenumber/\inserttotalframenumber
}
\setbeamercolor{footline}{fg=blue}
\setbeamerfont{footline}{series=\bfseries}

% Let's get started
% BEGIN MAIN DOCUMENT
\begin{document}

\begin{frame}
  \titlepage
\end{frame}

\begin{frame}{ESPEC Web Controller, V3}
  \tableofcontents
  % You might wish to add the option [pausesections]
\end{frame}
%

% SECTION 1 %%%%%%%%%%%%%%%%%%%%%%%%%%%%%%%%%%%%%%%%%%%%%%%%%%%%%%%%
\section{1. Introduction: T-sereis \& TM, Q-link and EWC}

\begin{frame}{1.0: Introduction}{} 
What we will cover in this training session: 

\begin{itemize} 
    \item Brief overview of T-series offerings
    \item How and where does EWC come in? \\ Advantages? Benefits? 
    \item What is EWC V3? 
    \item EWC V3 overall functionality (but brief): \\
Hands-on Demonstrations
    \item Chamber Configuration \& Communication, Networking
    \item Troubleshooting: Chamber Configuration, Networking, and All That 
    \item U-Web Kit on existing T-series chambers
    \item Conclusion: Questions, comments, discussion.
\end{itemize} 
\end{frame}

\begin{frame}{1.1: HALT/HASS \& Extreme Temp Chambers}{} 
HALT/HASS chambers offered with Typhoon Manager (TM) or Q-link option.   
\begin{figure} [hbt!] 
\centering
\includegraphics[scale=0.18]{figures/t-series-and-TM-001.PNG}
\caption{HALT/HASS chambers with TM and/or Q-link option}
\end{figure}

\begin{itemize} 
    \item T-series with TM operates as a standalone system. 
    \item T-series with TM and Q-link option can be operated
          locally or remotely via Q-link Control Client. 
\end{itemize} 
\end{frame}

% new slide 

\begin{frame}{1.2: ESPEC Web Controller V3 for T-series}{} 
Customers can choose a T-series between TM, TM and Q-link or ESPEC Web Controller V3 (EWC).
\begin{figure} [hbt!]
\centering
\includegraphics[scale=0.29]{figures/Tseries-to-EWC-001.PNG}
\caption{EWC options/benefits over TM and/or Q-link}
\end{figure}
\end{frame} 

% end of new slide 

\begin{frame}{1.3: T-Series w/ EWC as Standalone System}{} 
T-series w/ EWC is offered with a touch display for its main operation as dedicated HMI.  

\begin{figure} [hbt!] 
\centering
\includegraphics[scale=0.14]{figures/tseries-ewc-touchscreen-001.PNG}
\caption{EWC with its touch display as a standalone system}
\end{figure}

\begin{itemize}
   \item Replaces the Control PC
   \item Touch display offers convenient operation
   \item Options: Standalone system or networked system

\end{itemize} 
\end{frame} 

\begin{frame}{1.4: T-Series w/ EWC as Networked System}{} 
T-series w/ EWC offers flexibility as a standalone or networked system.  
\begin{figure} [hbt!] 
\centering
\includegraphics[scale=0.22]{figures/T-series-ewc-networked-001.PNG}
\caption{EWC option as standalone or networked system}
\end{figure}
Networked System:
\begin{itemize}
   \item No dedicated Control PC is required.
   \item EWC can be accessed and operated remotely by authorized users from any device on the network. 
\end{itemize} 
\end{frame} 

% SECTION 2 %%%%%%%%%%%%%%%%%%%%%%%%%%%%%%%%%%%%%%%%%%%%%%%%%%%%%%%%
\section{2. TM vs EWC User Interface} 

\begin{frame}{2.0: TM Operation \& User Interface}{}
{\it{Standard}} offering: T-series system requires a Typhoon Manager (TM) software installed on a dedicated computer (called, Control PC) to operate the chamber. Control of chamber is done locally. Communication between T-series and Control PC is via TCP/IP. 
\begin{figure} [hbt!] 
\centering
\includegraphics[scale=0.26]{figures/Tseries-TM-UI-001.PNG}
\caption{T-Series with TM software on Control PC}
\end{figure}
\end{frame}            
%
\begin{frame}{2.1: EWC Operation \& User Interface}{}
T-series with EWC V3 eliminates the need of a Control PC (with TM). EWC 'replaces' TM per operational features specified in the chamber. With its touch display, EWC offers chamber control/operation locally (similar to TM) as a single unit. 
\begin{figure} [hbt!] 
\centering
\includegraphics[scale=0.26]{figures/Tseries-EWC-UI-001.PNG}
\caption{T-Series with EWC and its touch display}
\end{figure}
\end{frame}

\begin{frame}{2.2: EWC User Interface: Networked System}{}
T-series with EWC V3 can join a network to allow its operation and control via devices remotely. 

\begin{itemize}
   \item Operator can control the chamber locally via the touch display.
   \item Operators can access EWC to control the chamber remotely (for authorized users only). 
\end{itemize} 

\begin{figure} [hbt!] 
\centering
\includegraphics[scale=0.24]{figures/Tseries-and-EWC-network-001c.PNG}
\caption{EWC accessible by authorized users on a network}
\end{figure}
\end{frame}

% SECTION 3 %%%%%%%%%%%%%%%%%%%%%%%%%%%%%%%%%%%%%%%%%%%%%%%%%%%%%%%%
\section{3. What is ESPEC Web Controller (EWC), V3?} 

\begin{frame}{3.0: Introducing ESPEC Web Controller, V3}{} 
ESPEC Web Controller V3 (EWC) is an embedded computer.  

\begin{itemize}
   \item {\bf{Hardware}}: Based on an x86 architecture, called Up2Board (pictured below). 
   \item {\bf{Software}}: Powered by Debian GNU/Linux distribution as its operating system (OS). 
\end{itemize} 

\begin{figure} [hbt!] 
\centering
\includegraphics[scale=0.22]{figures/tseries-up2-001.png}
%\includegraphics[scale=0.22]{figures/tseries-up2-002.PNG}
\caption{EWC Up2Board, its I/O ports (discuss later)}
\end{figure}
\end{frame}

% new slide 

\begin{frame}{3.1: ESPEC Web Controller V3: Features}{} 
\begin{itemize}
   \item EWC V3 uses a standard Web browser to display (i.e., host) its user 
         interface (UI) for controlling the chamber; hence, the name {\bf{Web Controller}}.  

   \item EWC V3 can operate as a standalone machine with a touch or standard display. 
         Its UI is displayed using a subset of a Web browser, called Electron. 
    
   \item EWC V3 can join a network to host its UI through a Web browser, accessable 
         through its IP address or hostname. Thanks to the ability of a Web 
         browser to operate from any computer on the network, chamber operations 
         can be performed via EWC V3 remotely by authorized users on the network.  

   \item No device outside the network can access EWC V3, unless the network is 
         configured to host port forwarding to make EWC V3 visible on the Internet. 

\end{itemize}
\end{frame} 

% end of new slide 

\begin{frame}{3.2: Updates \& Security Features}{} 
\begin{itemize}
% Here is another example of security issues.
   \item {\bf{Security:}} With Debian GNU/Linux periodically receiving security updates, 
              EWC V3 is kept extremely secure. EWC V3 is configured with read-only 
              access to prevent intruders from exploiting its system. 

   \item {\bf{Updates \& Stability:}} EWC V3 firmware update is managed by a cloud 
              service provider (called MenderIO)  which employs a dual-root (stepladder) 
              technique to provide a smooth update. One root system is always 
              in a stable state to ensure an uninterrupted operation, should the firmware update fail.

   \item {\bf{Access Privilege \& User Policy:}} EWC V3 is shipped with one user account 
              (called {\bf{admin}}) with complete powers to manage the system. 
              This {\bf{admin}} account can enforce a user policy on EWC V3 and 
              administer additional accounts with various privileges. 
              Only authorized users can log into and control EWC V3.

\end{itemize}
\end{frame} 

\begin{frame}{3.3: Cloud Updating}{}
EWC V3 uses chamber S/N as its registration number to receive firmware update 
through our cloud service provider, MenderIO.
    \begin{figure} [hbt!] 
      \centering
%      \includegraphics[scale=.30]{figures/menderio-cloud-001.PNG}  
      \includegraphics[scale=0.320]{figures/mender-cloud-service-002.PNG}
      \caption{EWC subscribes to MenderIO cloud service}
    \end{figure}
\end{frame} 

\begin{frame}{3.4: Firmware Update: Online or Offline}{} 
           Customer has complete control over firmware update method: Offline or 
           Online through a cloud service. 
\begin{enumerate}
\item Using a stepladder technique, update package is pushed onto the next 
      available root partition (see diagram). 

    \begin{figure} [hbt!] 
      \centering
%      \includegraphics[scale=.30]{figures/menderio-cloud-001.PNG}  
      \includegraphics[scale=.60]{figures/stepladder-001.PNG}
      \caption{Stepladder technique of EWC V3 update}
    \end{figure}

\item If the update process is successful, the system will boot and run on this new firmware.   
\item If the update process failed, the system rolls back to the previous EWC V3 
      known to be in the stable state. 
\end{enumerate} 
\end{frame} 

\begin{frame}{3.5: Firmware Update Subscription}{} 
Firmware update is provided through a yearly update subscription. 
Customers purchase the subscription to receive the updates. 
\begin{itemize}
   \item {\bf{Online Update}}: \\
       For EWC V3 that connects to the Internet, MenderIO cloud updating service handles 
       the update process. Automatic update will take place when the chamber 
       (connected EWC V3) is in a Standby mode.   
   \item {\bf{Offline Update}}:   \\
       For EWC V3 that does not connect to the Internet, the owner is provided with a link 
       to access their account to download an update package to perform the update 
       procedure on their system. Refer to User's Manual for detail on Firmware Update procedure. 
\end{itemize} 
\end{frame} 

% SECTION 4
\section{4. EWC User Interface: Display Options}
\begin{frame}{4.0: Display Options \& User Interface}{}

%\begin{minipage}[c]{0.50\textwidth}
%\vspace{0.125in} 

\begin{itemize} 
   \item With touch display, EWC provides on-screen keyboard for control/operaton.  
\end{itemize} 
%\end{minipage}
%\hfill
%\begin{minipage}[c]{0.450\textwidth}

\begin{figure} [hbt!] 
      \centering
%      \includegraphics[scale=.15]{figures/onscreen-login-001a.PNG} 
      \includegraphics[scale=0.17]{figures/onscreen-login-001a.PNG}  
%            \includegraphics[width=\textwidth]{figures/onscreen-login-001a.PNG}
      \caption{EWC log-in screen on a touch display}
%      \captionof{figure}{EWC V3 log-in screen on a touchscreen}
\end{figure}
%\end{minipage} 
%\vspace{0.2in} 

\begin{itemize} 
    \item User Interface via a Web browser (remote access). 
\end{itemize} 

\begin{figure} [hbt!] 
      \centering
      \includegraphics[scale=.15]{figures/desktop-login-001.PNG}  
      \caption{EWC log-in screen on a Web browser}
\end{figure}
\end{frame} 

\begin{frame}{4.1: User Interface: EWC Home Page}{}
EWC user interface versus TM user interface.   
    \begin{figure} [hbt!] 
      \centering
      \includegraphics[scale=.15]{figures/touch-display-UI-002.PNG}  
      \caption{EWC V3 Overview (Home) page}
    \end{figure}

\begin{itemize}  
   \item {\bf{Menu bar:}} Access to all EWC's menus for chamber control/operation \& more.  
   \item {\bf{Status bar:}} Quick summary of chamber current status. 
   \item {\bf{Main Display:}} Detailed info of chamber operation.  
   \item {\bf{Touch Bar:}} Direct access \& control (next slide).    
\end{itemize} 
\end{frame}

\begin{frame}{4.2: Touch Bar: Direct Operation}{}
Touch display and its touch bar (top bar) operation.
    \begin{figure} [hbt!] 
      \centering
      \includegraphics[scale=.20]{figures/top-control-bar-001.PNG}
      \includegraphics[scale=.20]{figures/touch-display-bar-comp-001.PNG}  
      \caption{Touch display direct operation buttons}
    \end{figure}
\begin{enumerate}  
   \item Paging: backward, forward, refresh, direct access, etc.
   \item Screen capture operation  
   \item Network information (IP address)
   \item Removable USB drive (eject button)    
\end{enumerate} 
\end{frame}

\begin{frame}{4.3: EWC and Its Menu Bar}{}
\begin{minipage}[c]{0.50\textwidth}
All features of EWC V3 are accessable and controllable through the nine menus in the menu bar (icons); this menu remains fixed on the left. 

\vspace{0.125in} 
The {\bf{Constant}}, {\bf{Program}}, {\bf{Start/Stop}} and {\bf{Settings}} menus are generally linked to the user's privileges controlled by the EWC V3 administrator. 

\vspace{0.125in}
These menus are grayed out to those users without read/write privileges. 
\end{minipage}
\hfill
\begin{minipage}[c]{0.45\textwidth}
%    \begin{figure} [hbt!] 
%      \centering
%      \includegraphics[scale=.50]{figures/main-menu-001.PNG}  
      \includegraphics[width=0.90\textwidth]{figures/Tseries-menu-bar-001.PNG} 
%      \caption{Main menu of EWC V3}
      \captionof{figure}{EWC V3 main menu}
%    \end{figure}

%      \includegraphics[width=\textwidth]{figures/main-menu-001.PNG} 
%      \caption{Main menu of EWC V3}
%      \captionof{figure}{EWC V3 main menu}
\end{minipage}
\end{frame}

\begin{frame}{4.4: The Main Menu \& Features}{} 

\begin{enumerate}
   \item {\bf{User}}: log in/out; password reset or rescue.
   \item {\bf{Overview}}: Quick summary of chamber operating status/conditions.
   \item {\bf{Trend}}: Scatter plot of data points from tests.
   \item {\bf{History}}: Logs of chamber operation and its statistics.
   \item {\bf{Constant}}: Setting or re-setting new constant values.
   \item {\bf{Program}}: Construct program instructions, preview program output; program file manipulation, and more.
   \item {\bf{Start Stop}}: Control chamber operating mode.
   \item {\bf{Settings}}: Settings of various features and operations. \\ This menu is extensive and should only be accessible to (and controlled by) the administrator. 
   \item {\bf{About}}: Information about terms and agreement of EWC and its firmware version, online support and embedded operation manual (User's Manual).
\end{enumerate} 
\end{frame}

\begin{frame}{4.5: EWC Functionality: Operation of Menus}{}
We now examine briefly EWC menus and their functionality.  Hands-on and demos...
\vspace{0.125in}

\begin{minipage}[c]{0.550\textwidth}
{\bf{Program Status:}} 
{\bf{Overview}} displays program status, step being executed, remaining time, etc.   
\end{minipage}
\hfill
\begin{minipage}[c]{0.40\textwidth}
%      \begin{figure} [hbt!] 
%        \centering
%        \includegraphics[scale=.20]{figures/f4t-network-cfg-001.PNG}  
        %\includegraphics[width=\textwidth]{figures/up2board-io-002.PNG}
        \includegraphics[width=\textwidth]{figures/tseries-prog-run-001.PNG}  
%        \caption{Example of F4T internal network setting}
        \captionof{figure}{Program execution}
%      \end{figure}
\end{minipage}

\begin{minipage}[c]{0.550\textwidth}
{\bf{Scatter Plot:}} 
Data points collected in chamber are plotted against time. Plotting refreshes continusouly to reflect new data.
\end{minipage}
\hfill
\begin{minipage}[c]{0.40\textwidth}
%      \begin{figure} [hbt!] 
%        \centering
%        \includegraphics[scale=.20]{figures/f4t-network-cfg-001.PNG}  
        \includegraphics[width=\textwidth]{figures/tSeries-plot-001.PNG}  
        %\includegraphics[width=\textwidth]{figures/tSeries-plots-001.PNG} 
%        \caption{Example of F4T internal network setting}
        \captionof{figure}{Trend graph}
%      \end{figure}
\end{minipage}

\vspace{0.125in} 

\begin{minipage}[c]{0.550\textwidth}
{\bf{Program List:}} 
Programs on the list can be viewed, edited, save, downloaded or removed. Programs are stored on the Web Controller storage system.  
\end{minipage} 
\hfill
\begin{minipage}[c]{0.400\textwidth} 
        \includegraphics[width=\textwidth]{figures/Tseries-prog-list-001.PNG} 
        \captionof{figure}{Program list}
%      \end{figure}
\end{minipage}
\end{frame}


\begin{frame}{4.6: EWC Functionality: Operation of Menus}{}
The {\bf{Settings}} menu and its submenus functionality.   
\vspace{0.125in}

\begin{minipage}[c]{0.550\textwidth}
{\bf{Network:}} 
DHCP \& static network settings. Excellent place for troubleshooting networking issues. 
\end{minipage}
\hfill
\begin{minipage}[c]{0.40\textwidth}
%      \begin{figure} [hbt!] 
%        \centering
%        \includegraphics[scale=.20]{figures/f4t-network-cfg-001.PNG}  
        %\includegraphics[width=\textwidth]{figures/up2board-io-002.PNG}
        \includegraphics[width=0.80\textwidth]{figures/ewc-network-001.png}  
%        \caption{Example of F4T internal network setting}
        \captionof{figure}{Program execution}
%      \end{figure}
\end{minipage}

\begin{minipage}[c]{0.550\textwidth}
{\bf{Chamber Interface:}} 
Chamber configuration and optional features, communication and settings. 
\end{minipage}
\hfill
\begin{minipage}[c]{0.40\textwidth}
%      \begin{figure} [hbt!] 
%        \centering
%        \includegraphics[scale=.20]{figures/f4t-network-cfg-001.PNG}  
        \includegraphics[width=0.80\textwidth]{figures/ewc-interface-001.png}  
%        \caption{Example of F4T internal network setting}
        \captionof{figure}{Trend graph}
%      \end{figure}
\end{minipage}

\vspace{0.125in} 

\begin{minipage}[c]{0.550\textwidth}
{\bf{Services:}} 
Services and resources, where different services can be managed, such as restarting without having to reboot EWC. Simple network issue can be fixed here.  
\end{minipage} 
\hfill
\begin{minipage}[c]{0.400\textwidth} 
        \includegraphics[width=0.80\textwidth]{figures/ewc-service-001.png} 
        \captionof{figure}{Program list}
%      \end{figure}
\end{minipage}
\end{frame}

\begin{frame}{4.7: User Interface Settings}{}
User Interface Settings offers direct control of volume and alarm settings (Touch display \& Web browser display).    
    \begin{figure} [hbt!] 
      \centering
      \includegraphics[scale=.15]{figures/touchscreen-volume-test.png}  
      \caption{Volume setting and testing on touch display}
    \end{figure}
\begin{enumerate}  
   \item Alarm: Enable or disable built-in speaker.
   \item Volume: Set volume level. 
   \item Test Settings: Test current vloume setting.   
\end{enumerate} 
\end{frame}

\begin{frame}{4.8: Touch Display: HMI Settings}{}
Touch display HMI setting features.    
    \begin{figure} [hbt!] 
      \centering
      \includegraphics[scale=.22]{figures/hmi-settings-001.PNG}  
      \caption{HMI setting on touch display}
    \end{figure}
\begin{itemize}  
   \item Language and Time Zone settings
   \item Display Timeout: Set idle time for monitor sleep mode
   \item Automatic Login: set auto login without password
\end{itemize} 
\end{frame}

\begin{frame}{4.9: Shell Terminal Access \& Troubleshooting}{}
Shell terminal access allows direct text terminal operation, to be used for troubleshotting the system (for local operator).    
    \begin{figure} [hbt!] 
      \centering
      \includegraphics[scale=.20]{figures/shell-terminal-001.PNG}  
      \caption{Direct access to embedded system OS}
    \end{figure}
\end{frame}

\begin{frame}{4.10: Firmware Version and Online User Manual}{}
EWC Firmware version and online User's Manual can be found under the {\it{About}} menu.     
    \begin{figure} [hbt!] 
      \centering
      \includegraphics[scale=.20]{figures/ewc-user-manual-001.png}  
      \caption{Firmware version and User Manual}
    \end{figure}
\end{frame}

% SECTION 5%%%%%%%%%%%%%%%%%%%%%%%%%%%%%%%%%%%%%%%%%%%%%%%%%%%%%%%%%%%%
\section{5. Chamber Configuration \& Interface} 

\begin{frame}{5.0: Configuration \& Communication Interface}{}

EWC communicates with the T-series using TCP/IP protocol for simple interfacing but fast throughout.  

    \begin{figure} [hbt!] 
      \centering
      \includegraphics[scale=.20]{figures/tseries-up2-002.PNG}  
      \caption{EWC hardware Ethernet ports}
    \end{figure}
 
 \begin{itemize}  
   \item {\bf{Main Network}}: Connects to EPO (external Ethernet port) of T-series chamber allowing EWC to join customer's main network for multiple access points.
   \item {\bf{T-Series}}: Connects directly to T-series PLC using its default internal network protocol. 
\end{itemize} 
\end{frame}

\begin{frame}{5.1: TCP/IP: A Closer Look}{}
\begin{minipage}[c]{0.550\textwidth}
{\bf{Internal Network:}} 

T-Series chambers use an internal network protocol with unique IP address to
communicate with EWC V3. This internal network is configured by EWC V3 by default, as illustrated in the figure (right).    
\end{minipage} 
\hfill
\begin{minipage}[c]{0.400\textwidth} 
        \includegraphics[width=\textwidth]{figures/tseries-tcp-01.PNG} 
%        \caption{Example of F4T internal network setting}
        \captionof{figure}{T-series network}
%      \end{figure}
\end{minipage}

\begin{minipage}[c]{0.550\textwidth}

{\bf{External Network:}} 

EWC V3 uses its predefined Ethernet port to connect to customer's main network to allow remote access for control/operation of the chamber.  
\end{minipage}
\hfill
\begin{minipage}[c]{0.40\textwidth}
%      \begin{figure} [hbt!] 
%        \centering
%        \includegraphics[scale=.20]{figures/f4t-network-cfg-001.PNG}  
        %\includegraphics[width=\textwidth]{figures/up2board-io-002.PNG}
        \includegraphics[width=\textwidth]{figures/UWeb-kit-001.png}  
%        \caption{Example of F4T internal network setting}
        \captionof{figure}{EWC network ports}
%      \end{figure}
\end{minipage}
\end{frame}

\begin{frame}{5.2: TCP/IP Options for T-series}{}

For T-Series chambers, internal network between EWC V3 and PLC is available only for LAN protocol via Network Interface options (see Figure). 

\vspace{0.125in} 

\begin{minipage}[c]{0.550\textwidth}

Default IP configurations are as follows: 

  \begin{itemize}
     \item {\bf{Network Interface}}: LAN
     \item {\bf{Network}}: Internal Network protocol: 121.165.141.0/24. 
     \item {\bf{Server IP}}: EWC V3 IP address: 121.165.141.155.
     \item {\bf{Server IP}}: PLC IP address: 121.165.141.61.
  \end{itemize} 
    
\end{minipage} 
\hfill
\begin{minipage}[c]{0.40\textwidth} 
        \includegraphics[width=\textwidth]{figures/tseries-tcp-01.PNG} 
%        \caption{Example of F4T internal network setting}
        \captionof{figure}{T-series LAN config}
%      \end{figure}
\end{minipage}
\end{frame}

\begin{frame}{5.3: Chamber Interface Config: Setup Wizard}{} 
Communication between EWC and T-series chamber is possible only when chamber interface configuration is done correctly, as shown below (factory setup).  

    \begin{figure} [hbt!] 
      \centering
      \includegraphics[scale=.230]{figures/t-series-cfg-001.PNG}  \\
      \caption{Factory Config: Setup Wizard chamber interface configuration}
    \end{figure}
\end{frame}


\begin{frame}{5.4: Chamber Interface Config: Settings Menu}{} 
Communication between EWC and T-series chamber can be re-established using the Settings menu to properly connect and configure the chamber, as shown below. 

    \begin{figure} [hbt!] 
      \centering
      \includegraphics[scale=.200]{figures/chamber-config-settings-001.png}  \\
      \caption{Settings menu: chamber interface configuration}
    \end{figure}
\end{frame}

% SECTION 6 %%%%%%%%%%%%%%%%%%%%%%%%%%%%%%%%%%%%%%%%%%%%%%%%%%%%%%%%%%%%
\section{6. Troubleshooting: Chamber Config. \& Networking, etc.}

\begin{frame}{6.0: Troubleshooting Communication Error}{} 
Example of chamber interface configuration error or no communication: chamber/controller does not respond to Web Controller. 

    \begin{figure} [hbt!] 
      \centering
      \includegraphics[scale=.220]{figures/web-controller-comm-issue-001.PNG}  \\
      \caption{EWC V3 attempting to communicate with the chamber}
    \end{figure}
Possible causes: 
\begin{itemize}
   \item Ethernet cable unplugged.
   \item Chamber is off.  
\end{itemize}
\end{frame}

\begin{frame}{6.1: No response from Chamber}{} 
Example of chamber interface configuration error; no communication: chamber/controller does not respond to Web Controller. After several attempts, Web Controller displays error message, 

    \begin{figure} [hbt!] 
      \centering
      \includegraphics[scale=.220]{figures/web-controller-comm-issue-002.PNG}  \\
      \caption{EWC V3 attempting to communicate with the chamber}
    \end{figure}
\end{frame}

\begin{frame}{6.2: Chamber Interface Config: Settings Menu}{} 
{\bf{Troubleshooting Procedure:}} \\ Chamber interface configuration under the {\bf{Settings}} menu, 
accessible via these buttons: {\keys{Settings}} $\rightarrow$ {\keys{Chamber Interface}}.   

\begin{itemize}
   \item {\bf{Chamber Category:}} T-Series Chamber. 
   \item {\bf{Model Selection:}} Select appropriate T-series chamber from T1.5 - T8.0 
   \item {\bf{Vibration:}} Optional Selection (based on available feature) 
   \item {\bf{Chamber COMM Interface:}} Internal LAN is the default network communication using TCP/IP protocol with a predefined IP address for EWC and T-series, all internally and automatically configured by EWC.
   \item {\bf{Optional Features:}} These features can be selected based on chamber's available options. 
\end{itemize} 
\end{frame}

\begin{frame}{6.3: Example of Chamber Configuration}{} 
Chamber configuration for T-series on EWC is quite straightforward, required only the correct type/model (T1.5, T2.0, etc) and available optional features in the {\bf{Optional Features}} column. Communication protocol is configured automatically by EWC.   

    \begin{figure} [hbt!] 
      \centering
      \includegraphics[scale=.240]{figures/tseries-cfg-005.PNG}  \\
      \caption{T-Series configuration: T2.5 w/ vibration}
    \end{figure}
\end{frame}

\begin{frame}{6.4: Troubleshooting Network Configuration}{} 
Troubleshoot static network setup for T-series w/ EWC. 
    \begin{figure} [hbt!] 
      \centering
      \includegraphics[scale=.24]{figures/TSeries-static.PNG}  \\
      \caption{Static config on T-series w/ EWC (touch display)}
    \end{figure}

\begin{itemize} 
   \item EWC uses Class C network with IP address: {\url{192.168.0.83}}
   \item Client PC must use the same subnet ({\url{255.255.255.0}}) with IP address {\url{192.168.0.xxx}}, e.g., {\url{192.168.0.84}}
   \item User opens Web browser on their remote PC and enter: {\url{http://192.168.0.83/}}.
\end{itemize} 
\end{frame}


\begin{frame}{6.5: Troubleshooting Network Configuration}{} 
For a DHCP network, EWC provides a smooth operation; simply connect and use.    
    \begin{figure} [hbt!] 
      \centering
      \includegraphics[scale=.24]{figures/TSeries-DHCP-001.PNG}  \\
      \caption{DHCP config on T-series w/ EWC (touch display)}
    \end{figure}

To access EWC remotely: 
\begin{enumerate} 
   \item Find out EWC IP address via touch bar (eth0) or Network settings submenu under Settings. 
   \item User opens Web browser on their remote PC and enter: {\url{http://IP-address/}} (from step 1).
\end{enumerate} 
\end{frame}

% SECTION 7 %%%%%%%%%%%%%%%%%%%%%%%%%%%%%%%%%%%%%%%%%%%%%%%%%%%%%%%%%%%%
\section{7. Universal Web Controller Kit (U-Web Kit) } 

\begin{frame}{7.0: U-Web Kit Support}{} 
Universal Web Controller kit (U-Web kit): Offers excellent solutions for control/operation of existing T-series chambers in place of or as an option to TM Control PC or Q-link. 
    \begin{figure} [hbt!] 
      \centering
      \includegraphics[scale=.02]{figures/up2board0a.PNG}  
      \caption{U-Web kit, EWC hardware}
    \end{figure}
{\bf{Advantages:}} 
\begin{itemize}
\item Portable and swappable between TM Control PC. 
\item Simple setup for standalone or networked operation. 
\item Customers have option for remote access.
\end{itemize} 
\end{frame}

\begin{frame}{7.1: U-Web Kit for T-series Chambers}{} 

Available Options: 
\begin{itemize}
    \item U-Web kit without touch display
    \item U-Web kit with touch display (?)
\end{itemize}  

    \begin{figure} [hbt!] 
      \centering
      \includegraphics[scale=.30]{figures/u-web-kit-001.PNG}  \\
      \caption{U-Web kit versus T-Series with EWC}
    \end{figure}
\end{frame}

\begin{frame}{7.2: U-Web Kit: Install Options}{} 

    \begin{figure} [hbt!] 
      \centering
      \includegraphics[scale=.35]{figures/Tseries-TM-EWC-U-Web-kit-001.PNG}  \\
      \caption{U-Web kit installation options}
    \end{figure}
\end{frame}


\begin{frame}{7.3: U-Web Kit: Network Setup (Static) }{} 
{\bf{Standalone system}}: Isolated from the rest of devices 
    \begin{figure} [hbt!] 
      \centering
      \includegraphics[scale=.26]{figures/U-Web-kit-static-cfg.PNG}  \\
      \caption{U-Web kit static network configuration}
    \end{figure}
Static config: 
\begin{itemize} 
   \item EWC uses Class C network with IP address: {\url{192.168.0.83}}
   \item Client PC must use the same subnet ({\url{255.255.255.0}}) with IP address {\url{192.168.0.xxx}}, e.g., {\url{192.168.0.84}}
\end{itemize} 
\end{frame}

\begin{frame}{7.4: U-Web Kit: Network Setup (DHCP)}{} 
{\bf{Networked system}}: Offers multiple remote access points for control/operation
    \begin{figure} [hbt!] 
      \centering
      \includegraphics[scale=.26]{figures/U-Web-kit-DHCP-cfg.PNG}  \\
      \caption{U-Web kit DHCP network configuration}
    \end{figure}
DHCP config: 
\begin{itemize} 
   \item EWC gets IP address from the DHCP server. 
   \item Client PC must be on the same network. 
\end{itemize} 
\end{frame}

\begin{frame}{7.5: U-Web Kit: Touch Display}{} 
U-Web kit options with touch display (open dicussion) 
\begin{itemize} 
   \item With touch display: 
         \begin{itemize} 
             \item Requires more work--essentially based on T-series with EWC manufactured at our factory.
             \item Not feasable to offer swabable option with existing TM Control PC. 
             \item Touch display is not really needed for this kit as U-Web kit's main purpose is to offer remote access or from a dedicated control PC. 
         \end{itemize}
   \item Without touch display: 
         \begin{itemize} 
             \item Installation is quick and easy, and can be done by customers; installation guide will be provided for Service Tech and DIY (Do-it-yourself) kit.
         \end{itemize} 
\end{itemize} 
\end{frame}


% SECTION 8 %%%%%%%%%%%%%%%%%%%%%%%%%%%%%%%%%%%%%%%%%%%%%%%%%%%%%%%%%%%%%%%%%
\section{8. Conclusions: Questions \& Answers} 

\begin{frame}{8.0: Conclusion}{}
EWC offers a number of advantages for T-series: 
\begin{itemize}
   \item Touch display for convenience and practility
   \item Offers as a standalone or networked system
   \item Simple connection and chamber configuration
   \item Easy troubleshooting

    \begin{figure} [hbt!] 
      \centering
      \includegraphics[scale=.120]{figures/prog-002.PNG}  
      \caption{Program/Preview menu of EWC V3}
    \end{figure}

   \item Questions \& Answers
\end{itemize}
\end{frame}


% All of the following is optional and typically not needed. 
\appendix
%\section<presentation>*{\appendixname}
%\subsection<presentation>*{References}
\section<presentation>*{References}

\begin{frame}[allowframebreaks]
  \frametitle<presentation>{References}
    
  \begin{thebibliography}{10}
    
  \beamertemplatebookbibitems
  % Start with overview books.

  \bibitem{ESPECWebManual2022}
    ESPEC.~Web Controller V3.
    \newblock {\em User Manual: ESPEC Web Controller, V3}, wiki/bitbucket.com, 2022. \\
    {\url{https://bitbucket.org/especnorthamerica/especweb/wiki/Web_Controller_V3_Home}}

  \bibitem{ESPECWebManual2022}
    ESPEC.~Web Controller V3.
    \newblock {\em User Manual} [ESPEC Web Controller, V3]. Retrievable and viewable
    via EWC menu bar under {\sl{About}}, 2022. \\
    {\url{http://especSN.local/\#/about}}

  \bibitem{ESPECWebConfig2022}
    ESPEC.~Production Configuration.
    \newblock {\em ESPEC Web Controller V3, Factory Config}, ESPEC Web Controller V3, Internal Departmental Resource, 2022. \\
    {\url{https://bitbucket.org/especnorthamerica/espec-web-vue/wiki/Factory_Setup_Procedure}}
    %\newblock MCC, 2020.
  

%  \beamertemplatearticlebibitems
%  \bibitem{UbuntuUnity2017}
%     Unity desktop ended.
%     \newblock Failed to popualrize or market Unity as a solution to all Ubuntu products, Unity project will be ended in April 2017, said Canonical boss Mark Shuttleworth.  Retrieved from {\url{https://www.bbc.com/news/technology-39490848}}. 


  %  S.~Someone.
  %  \newblock On this and that.
  %  \newblock {\em Journal of This and That}, 2(1):50--100,
  %  2000.
  \end{thebibliography}
\end{frame}

\end{document}
